\section{超参追踪与日志化}

\subsection{Logbook 结构}

课程建议为每次实验维护一份 logbook，记录以下字段：

\begin{itemize}
    \item \textbf{Learner IDs}：Actor、Critic、Target network 的 LR；
    \item \textbf{Replay ratio}：每个 update 消耗的 transitions；
    \item \textbf{Entropy schedule}：初始值 + decay 速率；
    \item \textbf{Benchmark metrics}：KL、critic loss、reward mean。
\end{itemize}

\begin{importantbox}{为什么要 log mutation}
比较不同 seed 时，往往只有超参数略微偏差。记录完整 logbook 能让你快速复制 good run，也能快速排查 catastrophic divergence。
\end{importantbox}

\subsection{自动化表格}

\begin{table}[H]
\centering
\begin{tabular}{p{0.25\textwidth}p{0.25\textwidth}p{0.25\textwidth}}
\toprule
字段 & 推荐工具 & 目的 \\
\midrule
KL curve & TensorBoard / WandB & 监控是否突破 trust region \\
Entropy & WandB custom scalar & 追踪 exploration collapse \\
Critic loss & CSV + daily summary & 发现 critic drift 前兆 \\
Replay ratio & elastic log & 确保 sample efficiency 不被忘记 \\
\bottomrule
\end{tabular}
\caption{Hyperparameter logging 的具体实践}
\end{table}

\subsection{本章小结}

超参日志化将实验从“跑几个 seed”升级为“可复现的工程流程”，对齐团队协作与 performance regression 检查。

